\subsection{指数记法}

设 M 为一个幺半群，以乘法记写，且 \(\mathbf{u} = (u_{x})_{x \in \mathbf{X}}\) 是 M 的一个元素族，两两交换。设 \(\alpha\) 属于 \(\mathbf{N}^{(\mathbf{X})}\) ；元素 \(u_{x}^{\alpha(x)}\) 和 \(u_{y}^{\alpha(y)}\) （M 中的）在 x, y 属于 X 时交换，并且存在 X 的一个有限子集 S，使得 \(u_{x}^{\alpha(x)} = 1\) 对 x 属于 X - S 成立。因此我们可以写出：

\[
\mathbf {u} ^ {\alpha} = \prod_ {x \in \mathrm{X}} u _ {x} ^ {\alpha (x)}.\tag{16}
\]

设 \(\mathbf{M}'\) 为 \(\mathbf{M}\) 的由族 \((u_x)_{x\in \mathbf{X}}\) 生成的子幺半群；它是交换的（§ 1，第 5 号，命题 4 的推论 2）。因此存在（第 7 号，命题 10）一个唯一的同态 \(f\) ，它把 \(\mathbf{N}^{(\mathbf{X})}\) 映入 \(\mathbf{M}'\) ，使得 \(f(\delta_x) = u_x\) 对所有 \(x\in \mathbf{X}\) 成立，且 \(f(\alpha) = \mathbf{u}^{\alpha}\) 对所有 \(\alpha\) （属于 \(\mathbf{N}^{(\mathbf{X})}\) ）成立。我们推导出以下公式

\[
\mathbf {u} ^ {\alpha + \beta} = \mathbf {u} ^ {\alpha}. \mathbf {u} ^ {\beta}\tag{17}
\]

\[
\mathbf {u} ^ {0} = 1\tag{18}
\]

\[
\mathbf {u} ^ {\delta_ {x}} = u _ {x}\tag{19}
\]

对 \(\alpha, \beta\) （在 \(\mathbf{N}^{(\mathbf{x})}\) 中）和 \(x\) （在 \(\mathbf{X}\) 中）成立。

设 \(\mathbf{v} = (v_x)_{x \in \mathbf{X}}\) 是 \(\mathbf{M}\) 的另一个元素族；假设 \(v_x v_y = v_y v_x\) 和 \(u_x v_y = v_y u_x\) 对 \(x, y\) （属于 \(\mathbf{X}\) ）成立。那么存在（§ 1，第 5 号，命题 4 的推论 2）一个交换子幺半群 \(\mathbf{L}\) （\(\mathbf{M}\) 的），使得 \(u_x \in \mathbf{L}\) 和 \(v_x \in \mathbf{L}\) 对所有 \(x \in \mathbf{X}\) 成立。映射 \(\alpha \mapsto u^\alpha . v^\alpha\) （从 \(\mathbf{N}^{(\mathbf{X})}\) 到 \(\mathbf{L}\) 中）则是一个同态（§ 1，第 5 号，命题 5），它把 \(\delta_x\) 映到 \(u_x . v_x\) 。因此我们有公式

\[
\mathbf {u} ^ {\alpha}. \mathbf {v} ^ {\alpha} = (\mathbf {u}. \mathbf {v}) ^ {\alpha},\tag{20}
\]

（其中 \(\mathbf{u}.\mathbf{v}\) 是族 \((u_x.v_x)_{x\in \mathbf{X}}\)

当 M 是交换的时， \(u^{\alpha}\) 可对 M 的元素的每个族 u 定义，且公式 (15) 至 (20) 无条件成立。

